\documentclass[a4,10pt]{aleph-notas}

% -- Paquetes adicionales
\usepackage{enumitem}
\usepackage{aleph-comandos}

% -- Datos
% El siguiente texto presenta la definición de la jerarquía acumulativa de von Neumann
% y demuestra que sus niveles cubren todos los conjuntos (la clase universal).
% Adaptación y exposición didáctica basada en el tratamiento estándar (ver Corolario 1.3.6).
\institucion{Proyecto Alephsub0}
\asignatura{Matemáticas de la Información}
\tema{Jerarquía de von Neumann: definición y universalidad}
\autor{Andrés Merino}
\fecha{Abril 2026}

% Fuente tipográfica y logotipo
\fuente{montserrat}
\logouno[4.5cm]{Logos/LogoAlephsub0-02}
\definecolor{colordef}{cmyk}{0.81,0.62,0.00,0.22}

% -- Comandos adicionales
\newtheorem*{prob}{Problema}
    \tcolorboxenvironment{prob}{%
        color=colordef,recuadrost,colback=colordef!7,drop fuzzy shadow
    }


\begin{document}


\encabezado

\noindent
% El siguiente problema corresponde al Ejemplo 1.3.7 del libro Introduction to Cardinal Aritmetic de M. Holz. Primero, presentemos algunas definiciones:

\begin{prob}
Sea $V=\mathbb{C}^2$ con la base
\[
    \mathcal{B}=\left\{
        b_1=\begin{pmatrix}1\\ i\end{pmatrix},
        b_2=\begin{pmatrix}1\\ -i\end{pmatrix}
    \right\}.
\]
\begin{enumerate}
    \item Encontrar la base dual $\mathcal{B}^*$.
    \item Expresar el funcional $f(x,y)=x+2y$ como combinación lineal de $\mathcal{B}^*$.
    \item Calcular
    \[
        b_1^*\begin{pmatrix}2+i\\ 1-i\end{pmatrix}
        \quad \text{y} \quad
        b_2^*\begin{pmatrix}2+i\\ 1-i\end{pmatrix}.
    \]
\end{enumerate}
\end{prob}

\begin{proof}
Sea $\mathcal{B}^*=\{b_1^*,b_2^*\}$ la base dual de $\mathcal{B}$. Por definición,
\[
    b_i^*(b_j)=\delta_{ij}.
\]

Para encontrar $b_1^*$, escribimos
\[
    b_1^*(x,y)=\alpha x+\beta y.
\]
Como $b_1^*(b_1)=1$ y $b_1^*(b_2)=0$, se tiene
\[
    \alpha+i\beta=1,
    \qquad
    \alpha-i\beta=0.
\]
Resolviendo el sistema,
\[
    \alpha=\frac{1}{2},
    \qquad
    \beta=-\frac{i}{2}.
\]
Por tanto,
\[
    b_1^*(x,y)=\frac{1}{2}x-\frac{i}{2}y.
\]

De forma análoga, para $b_2^*$ escribimos
\[
    b_2^*(x,y)=\alpha x+\beta y.
\]
Como $b_2^*(b_1)=0$ y $b_2^*(b_2)=1$, se tiene
\[
    \alpha+i\beta=0,
    \qquad
    \alpha-i\beta=1.
\]
Resolviendo el sistema,
\[
    \alpha=\frac{1}{2},
    \qquad
    \beta=\frac{i}{2}.
\]
Así,
\[
    b_2^*(x,y)=\frac{1}{2}x+\frac{i}{2}y.
\]
Por tanto,
\[
    \mathcal{B}^*
    =
    \left\{
        b_1^*(x,y)=\frac{1}{2}x-\frac{i}{2}y,
        \quad
        b_2^*(x,y)=\frac{1}{2}x+\frac{i}{2}y
    \right\}.
\]

Ahora expresamos $f(x,y)=x+2y$ como combinación lineal de la base dual:
\[
    f=\lambda b_1^*+\mu b_2^*.
\]
Como los coeficientes en la base dual se obtienen evaluando $f$ en los vectores de la base original,
\[
    \lambda=f(b_1)=f(1,i)=1+2i,
\]
y
\[
    \mu=f(b_2)=f(1,-i)=1-2i.
\]
Por tanto,
\[
    f=(1+2i)b_1^*+(1-2i)b_2^*.
\]

Finalmente,
\[
    b_1^*\begin{pmatrix}2+i\\ 1-i\end{pmatrix}
    =
    \frac{1}{2}(2+i)-\frac{i}{2}(1-i)
    =
    \frac{1}{2},
\]
y
\[
    b_2^*\begin{pmatrix}2+i\\ 1-i\end{pmatrix}
    =
    \frac{1}{2}(2+i)+\frac{i}{2}(1-i)
    =
    \frac{3}{2}+i.
\]
\end{proof}



\end{document}